% The architectural break, and the handoff in both directions.
\begin{tikzpicture}[font=\sffamily\scriptsize, x=1cm, y=1cm]

\node[chlabel, anchor=west] at (-6.8,3.6) {the same node, two completely different architectures, twenty milliseconds apart};

% ================================================================ the loader
\node[regcell, fill=hwcol, minimum width=54mm, text width=52mm, inner sep=3pt,
      anchor=north west, align=left] at (-6.8,3.2)
  {\textbf{the loader}\\[2pt]
   a polled superloop\\
   \textbf{interrupts disabled}\\
   two accepted identifiers, set in the hardware filter\\
   no scheduler, no tasks, no queues\\
   one flash driver and one checksum};

\node[regcell, fill=kercol, minimum width=54mm, text width=52mm, inner sep=3pt,
      anchor=north west, align=left] at (0.2,3.2)
  {\textbf{the application}\\[2pt]
   a task set on a kernel\\
   interrupts on, priorities assigned\\
   everything chapters 1 to 18 built\\
   a control task at one kilohertz\\
   middleware, safety, the bus, the loop};

% ================================================================ why polled
\node[chlabel, anchor=west] at (-6.8,0.6) {and why the left-hand column is polled, which is not laziness};
\node[regcell, fill=white, minimum width=114mm, text width=112mm, inner sep=3pt,
      anchor=north west, align=left] at (-6.8,0.2)
  {\textbf{1.} Interrupts remove the determinism a loader needs.\hfill
   \textbf{2.} Most parts cannot execute handlers during flash programming
   anyway.\\[2pt]
   \textbf{3.} The acceptance filter admits exactly two identifiers, so there is
   almost nothing to be interrupted about. Any one of the three is sufficient on
   its own, and all three come from primary documents rather than from
   preference.};

% ================================================================ the handoff
\node[chlabel, anchor=west] at (-6.8,-1.8) {the handoff, in both directions, through one RAM word};

\node[fw, text width=26mm, minimum height=9mm] (ld) at (-4.6,-2.6) {the loader};
\node[hw, text width=28mm, minimum height=9mm] (ram) at (-0.8,-2.6)
  {one RAM word\\\textbf{not initialised at startup}};
\node[laterblk, text width=26mm, minimum height=9mm] (ap) at (3.0,-2.6) {the application};
\draw[flow] (ld) -- (ram);
\draw[flow] (ram) -- (ap);
\draw[recoveredge] (ap) to[bend left=18] node[buslabel]{set the keyword, then reset} (ld);

\node[regcell, fill=red!18, minimum width=114mm, text width=112mm, inner sep=3pt,
      anchor=north west, align=left] at (-6.8,-3.9)
  {\textbf{Why not simply jump?} Because a direct jump leaves peripherals
   initialised by the loader, and an application developed and tested without a
   loader present, in the words of a primary source, \textbf{"can fail for subtle
   reasons that can be very difficult to track down"}. Those reasons surface in
   the field rather than on the bench, which is the worst place to find them and
   the reason this word exists.};

\node[umlnote, text width=114mm, anchor=north west] at (-6.8,-5.9)
  {For a volume that has spent eighteen chapters on a task set, the loader is a
   genuine architectural discontinuity, and it is worth saying so out loud rather
   than letting a reader assume the two halves are built the same way. The two
   columns above share a part, a bus and a flash driver, and share nothing else
   at all.};
\end{tikzpicture}
